\documentclass[12pt]{extarticle}
\linespread{1.5}

\title{ALU\_Register}
\author{Владимир Попов Сергеевич, Б01-411, Вариант 1, popov.vs@phystech.edu}
\date{\today}

\usepackage{cmap}
\usepackage[T2A]{fontenc}
\usepackage[english, russian]{babel}
\usepackage[utf8]{inputenc}
\usepackage{mathtools}
\usepackage{amsfonts}
\usepackage{graphicx}
\usepackage{listings}
\usepackage{courier}
\usepackage{indentfirst}
\usepackage{stmaryrd}
\usepackage{caption}
\usepackage{MnSymbol,wasysym}
\captionsetup[figure]{labelformat=empty}

\usepackage{tikz}
\usepackage{tzplot}
\usetikzlibrary{arrows,decorations.pathmorphing,backgrounds,positioning,fit,petri}
\usetikzlibrary{calc,intersections,through,backgrounds, quotes, angles}
\usetikzlibrary{arrows,automata,positioning}
\tikzset{snake it/.style={decorate, decoration=snake}}
\usepackage{pgfplots,pgf-pie}

\usepackage{geometry} % Меняем поля страницы
\geometry{left=2cm}% левое поле
\geometry{right=1.5cm}% правое поле
\geometry{top=1cm}% верхнее поле
\geometry{bottom=2cm}% нижнее поле

\newtheorem{definition}{Определение}
\newtheorem{remark}{Ремарка!}
\newtheorem{theorem}{Теорема}

\newcommand{\implyarrow}{%
  \mathrel{\raisebox{1.3ex}{\rotatebox[origin=c]{90}{\mathhexbox37F}}}}

\DeclareMathOperator{\arccot}{arccot}

\usepackage{listings}
\usepackage{xcolor}

%New colors defined below
\definecolor{codegreen}{rgb}{0,0.6,0}
\definecolor{codegray}{rgb}{0.5,0.5,0.5}
\definecolor{codepurple}{rgb}{0.758,0,0.82}
\definecolor{backcolour}{rgb}{0.95,0.95,0.92}

%Code listing style named "mystyle"
\lstdefinestyle{mystyle}{
  backgroundcolor=\color{backcolour},   commentstyle=\color{codegreen},
  keywordstyle=\color{magenta},
  numberstyle=\tiny\color{codegray},
  stringstyle=\color{codepurple},
  basicstyle=\ttfamily\footnotesize,
  breakatwhitespace=false,         
  breaklines=true,                 
  captionpos=b,                    
  keepspaces=true,                 
  numbers=left,                    
  numbersep=5pt,                  
  showspaces=false,                
  showstringspaces=false,
  showtabs=false,                  
  tabsize=2
}

%"mystyle" code listing set
\lstset{style=mystyle}

\begin{document}

\maketitle

\newpage

\tableofcontents

\newpage

\section{Функциональное описание}

Разработанный модуль АЛУ выполняет определённые входным кодом арифметические операции над двумя операндами и подаёт на выход результат. Данный АЛУ тактируется положительным фронтом, то есть меняет состояние внутреннего регистра и выдаёт следующий результат по положительному фронту входного сигнала \texttt{clk\_i}. Также этот АЛУ имеет возможность синхронного сброса, то есть, если при положительном фронте сигнала \texttt{clk\_i} входной сигнал \texttt{rst\_i=1}, тогда значение регистра и, соответственно, выходного сигнала сбрасываются в 0. В таблице \ref{tab:opcodes} представлены какому коду операции соответствует какое действие, производимое АЛУ.

\begin{table}[h!]
    \centering
    \caption{Коды операций разработанного АЛУ}
    \label{tab:opcodes}
    \begin{tabular}{|c|l|}
        \hline
        \textbf{opcode\_i}  & \textbf{Операция над first\_i и second\_i} \\ \hline
        \texttt{3'b000}     & Побитовое НЕ-И \\ \hline
        \texttt{3'b001}     & Побитовое Исключающее ИЛИ \\ \hline
        \texttt{3'b010}     & Сложение (без знака) \\ \hline
        \texttt{3'b011}     & Арифметический (знаковый) сдвиг \texttt{first\_i} вправо на \texttt{second\_i} \\ \hline
        \texttt{3'b100}     & Побитовое ИЛИ \\ \hline
        \texttt{3'b101}     & Логический сдвиг \texttt{first\_i} влево на \texttt{second\_i} \\ \hline
        \texttt{3'b110}     & Побитовое НЕ для \texttt{first\_i} \\ \hline
        \texttt{3'b111}     & \texttt{first\_i} $<$ \texttt{second\_i} \\ \hline
    \end{tabular}
\end{table}

\section{Структурная схема}

На стрелках, ведущими ко входам мультиплексора указаны индексы, по которым выбирается этот вариант. Все логические вентили представленные на схеме являются побитовыми

\begin{figure}[h!]
    \centering
    \includegraphics[width=1\linewidth]{diagram.png}
    \caption{Диаграмма ALU}
    \label{fig:diagram}
\end{figure}

\newpage

\section{Параметры}

\begin{table}[h!]
    \centering
    \caption{Описание параметров}
    \label{tab:params}
    \begin{tabular}{|c|c|l|}
        \hline
        \textbf{Название}  & \textbf{Значениe}                  & \textbf{Описание} \\ \hline
        \texttt{WIDTH}     & Размер слова (по умолчанию 64)     & Количество бит в операндах \\ \hline
    \end{tabular}
\end{table}

\newpage

\section{Порты}

\begin{table}[h!]
    \centering
    \caption{Описание портов}
    \label{tab:ports}
    \begin{tabular}{|c|c|l|}
        \hline
        \textbf{Название}  & \textbf{Ширина}  & \textbf{Описание} \\ \hline
        \texttt{clk\_i}     & 1       & Тактовый сигнал \\ \hline
        \texttt{rst\_i}     & 1       & Сигнал сброса \\ \hline
        \texttt{opcode\_i}  & 3       & Код выполняемой операции \\ \hline
        \texttt{first\_i}   & WIDTH   & Первый операнд \\ \hline
        \texttt{second\_i}  & WIDTH   & Второй операнд \\ \hline
        \texttt{result\_o}  & WIDTH   & Результат операции \\ \hline
    \end{tabular}
\end{table}

\section{Тактирование и сброс}

Тактирование осуществляется через входной порт clk\_i. Запись в D-триггер производится по положительному фронту. 

Сброс осуществляется через входной порт rst\_i. Сброс синхронный с тактовым сигналом, то есть по положительному фронту clk\_i в D-триггер записывается вычисленное значение, если rst\_i=0, и записывается 0, если rst\_i=1. Сброс в 0 происходит при положительном фронте тактового сигнала и rst\_i выставленным в 1.

\section{Тестирование}

В always-блоке описывается изменение тактового сигнала с 0 на 1 каждую единицу времени. Далее будем называть циклом - время между 2 ближайшими положительными фронтами тактового сигнала.

В начале обнулим rst\_i, после синхронизируемся с тактовым сигналом, выставим rst\_i в единицу, чтобы сбросить модуль в начальное состояние, опять синхронизируемся с тактовым сигналом, вернём rst\_i в ноль. 

Далее будем по положительному фронту тактового сигнала выставлять поочерёдно все возможные коды операций в opcode\_i, а также рандомить первый и второй операнд. Рандомить мы будем с конкретным сидом, чтобы можно было воспроизвести тестовый набор. Так, чтобы написать тесты, можно для каждого сида руками посчитать корректные значения и сравнивать с ними.

\end{document}